% =============================================================================
%  Capítulo 1 — Introducción (se ESCRIBE AL FINAL, alto nivel).
% =============================================================================
\chapter{Introducción}
\label{cap:introduccion}

\section{Alcance}
\label{sec:alcance}

El presente proyecto tiene como objetivo fundamental la construcción de un
sistema de Inteligencia Artificial generativa, altamente configurable, capaz
de sintetizar imágenes tanto en blanco y negro como en color.

Para lograr este objetivo, el alcance del sistema desarrollado abarca la
implementación de los siguientes componentes técnicos:

\begin{itemize}
  \item \textbf{Procesos de difusión para la inyección de ruido:}
        modelos \emph{Variance Exploding} (VE) y \emph{Variance Preserving} (VP).
  \item \textbf{Muestreadores (samplers):} integrador estocástico
        Euler--Maruyama, esquema Predictor--Corrector y \emph{Probability
        Flow ODE} determinista.
  \item \textbf{Programaciones de ruido (noise schedules):} decaimiento
        lineal y de tipo coseno.
  \item \textbf{Métricas de evaluación:} \emph{Negative Log-Likelihood} (NLL),
        \emph{Bits Per Dimension} (BPD), \emph{Fréchet Inception Distance}
        (FID) e \emph{Inception Score} (IS).
  \item \textbf{Generación controlable:} muestreo condicionado a clases
        (CFG) e imputación de imágenes parcialmente observadas.
\end{itemize}

\section{Estructura del documento}
\label{sec:estructura}

\todo{Reescribir cuando esté la versión final. Texto de partida:}

En el \textbf{Capítulo~\ref{cap:estado-arte}} se presenta el Estado del Arte,
con la formulación matemática de los modelos de difusión, los samplers, las
\emph{noise schedules}, las métricas de calidad y el problema de la
generación condicionada.

El \textbf{Capítulo~\ref{cap:desarrollo}} aborda el Desarrollo de Software y
Diseño, detallando el análisis de requisitos, los casos de uso, la
arquitectura del paquete \texttt{diffusion\_lib} y la planificación temporal.

En el \textbf{Capítulo~\ref{cap:resultados}} se exponen los Resultados,
evaluando el rendimiento, la robustez y la calidad gráfica del sistema
mediante comparativas y métricas cuantitativas.

Finalmente, el \textbf{Capítulo~\ref{cap:conclusiones}} recoge las
Conclusiones del trabajo, y el documento se cierra con la bibliografía y los
apéndices técnicos.
